% DO NOT EDIT. Generated from _data/cv_source.yml by scripts/build_cv.py.
% Edit cv/preamble.tex to change styling.

% Styling for the generated CV. Safe to edit by hand.
\documentclass[11pt,letterpaper,sans]{moderncv}

\moderncvstyle{classic}
\moderncvcolor{blue}

\usepackage[utf8]{inputenc}
\usepackage[scale=0.82]{geometry}
\setlength{\hintscolumnwidth}{3.2cm}
\setlength{\makecvtitlenamewidth}{10cm}

\usepackage{xurl}
\usepackage{enumitem}
\setlist[itemize]{leftmargin=*, itemsep=0.15em, topsep=0.15em}


\firstname{Lauren}
\familyname{Hoffman}
\email{lauren.hoffman@ucsc.edu}

\begin{document}

\makecvtitle
\vspace{-1.0em}

\section{Research Interests}
\cvitem{}{Polar science; sea ice; machine learning and explainable artificial intelligence; physical oceanography; air--sea--ice interactions; climate variability and predictability; AI-enabled approaches for polar observing applications.}

\section{Education}
\cventry{Sep. 2023}{Ph.D., Chemical Engineering}{University of California San Diego}{San Diego, CA}{}{}
\cventry{Jun. 2019}{M.S., Chemical Engineering}{University of California San Diego}{San Diego, CA}{}{}
\cventry{May 2016}{B.A./B.S., Mechanical Engineering}{University of San Diego}{San Diego, CA}{}{Mathematics and Chemistry minors; Honors Program. 3.98 major GPA; Dean's List, First Honors.}

\section{Research Experience}
\cventry{Jan. 2026--Present}{Postdoctoral Researcher}{University of California Santa Cruz}{Santa Cruz, CA}{PI: Xavier Prochaska}{%
\begin{itemize}
    \item Characterizing density-based ocean fronts and associated submesoscale dynamics in high-resolution ocean model output.
    \item Developing scalable analysis and visualization workflows for large ocean model datasets.
    \item Contributing machine learning, modeling, and data-analysis expertise to emerging Antarctic sea ice observing-system design efforts.
\end{itemize}}

\cventry{Feb. 2024--Dec. 2025}{Postdoctoral Researcher}{Université catholique de Louvain}{Louvain-la-Neuve, Belgium}{PI: François Massonnet}{%
\begin{itemize}
    \item Developed machine learning and explainable AI approaches to predict and understand rapid Arctic sea ice loss events on seasonal to interannual timescales.
    \item Applied interpretable neural networks and climate model ensembles to investigate drivers and predictability of sea ice variability.
\end{itemize}}

\cventry{Oct. 2023--Jan. 2024}{Postdoctoral Researcher}{Scripps Institution of Oceanography, UC San Diego}{San Diego, CA}{PI: Matt Mazloff}{%
\begin{itemize}
    \item Evaluated explainable machine learning techniques for geoscience applications.
    \item Implemented components of an operational forecast system for sea-ice motion prediction using machine learning.
\end{itemize}}

\cventry{Jun. 2019--Sep. 2023}{Graduate Researcher}{Scripps Institution of Oceanography, UC San Diego}{San Diego, CA}{PI: Matt Mazloff}{%
\begin{itemize}
    \item Investigated upper-ocean salinity response to atmospheric river events in the California Current System using observations and ocean model output.
    \item Developed machine learning and explainable AI approaches to predict and understand Arctic sea-ice dynamics from remote sensing data.
\end{itemize}}

\cventry{Sep. 2018--Mar. 2019}{Graduate Researcher}{University of California San Diego}{San Diego, CA}{PI: Andrea Tao}{%
\begin{itemize}
    \item Studied methods for improving recycled polymer blends using colloidal dispersions of silver nanoparticles.
\end{itemize}}

\cventry{Mar. 2018--Aug. 2018}{Research Assistant}{University of San Diego}{San Diego, CA}{PI: Truc Ngo}{%
\begin{itemize}
    \item Characterized 3D printed PMMA impregnated with drugs using supercritical carbon dioxide processing techniques.
\end{itemize}}

\cventry{Mar. 2015--Jun. 2017}{Undergraduate Research Assistant and Honors Thesis}{University of San Diego}{San Diego, CA}{PI: Truc Ngo}{%
\begin{itemize}
    \item Worked with members of El Cercado, Dominican Republic, and a community liaison to design a solar thermal water heater that met community needs and supported local ownership of the technology.
\end{itemize}}


\section{Industry Experience}
\cventry{Jun. 2016--Jun. 2018}{Lead Engineer}{Primo Wind, Inc.}{San Diego, CA}{}{%
\begin{itemize}
    \item Implemented design solutions for a small wind turbine system that improved power output and structural stability and optimized sustainability of materials and manufacturing.
\end{itemize}}


\section{Technical Expertise}
\cvitem{Programming}{Python; MATLAB; UNIX/Linux; LaTeX; BibTeX.}
\cvitem{Software/Tools}{TensorFlow/Keras; Git; Claude.}
\cvitem{Methods}{Machine learning; neural networks; explainable artificial intelligence; predictive modeling; data analysis; remote sensing; numerical modeling; scientific computing.}
\cvitem{Models}{MITgcm; LLC4320; CESM; CMIP6.}
\cvitem{Languages}{French, intermediate; Spanish, intermediate.}

\section{Publications}
\subsection{First-Author Publications}
\cvitem{In progress}{Hoffman, L., and F. Massonnet. eXplainable AI reveals sea surface temperature patterns associated with slowdowns in the decline of Arctic September sea ice extent in CESM2-LE.}
\cvitem{2025}{Hoffman, L., F. Massonnet, and A. Sticker. Probabilistic forecasts of September Arctic sea ice extent at the interannual timescale with data-driven statistical models. \textit{Journal of Geophysical Research: Machine Learning and Computation}, 2, e2025JH000669. DOI: 10.1029/2025JH000669}
\cvitem{2025}{Hoffman, L., M. Mazloff, S.T. Gille, D. Giglio, and P. Heimbach. Evaluating the trustworthiness of explainable artificial intelligence (XAI) methods applied to regression predictions of Arctic sea-ice motion. \textit{Artificial Intelligence for the Earth Systems}, 4, e240027. DOI: 10.1175/AIES-D-24-0027.1}
\cvitem{2023}{Hoffman, L., M. Mazloff, S.T. Gille, D. Giglio, P. Heimbach, C. Bitz, and K. Matsuyoshi. Machine learning for daily forecasts of Arctic sea-ice motion: an attribution assessment of model predictive skill. \textit{Artificial Intelligence for the Earth Systems}, 2, 230004. DOI: 10.1175/AIES-D-23-0004.1}
\cvitem{2022}{Hoffman, L., M. Mazloff, S.T. Gille, D. Giglio, and A. Varadarajan. Ocean salinity response to atmospheric river precipitation events in the California Current System. \textit{Journal of Physical Oceanography}, 52, 1867--1885. DOI: 10.1175/JPO-D-21-0272.1}
\cvitem{2018}{Hoffman, L., and T.T. Ngo. Affordable solar thermal water heating solution for rural Dominican Republic. \textit{Renewable Energy}, 115, 1220--1230. DOI: 10.1016/j.renene.2017.09.046}
\subsection{Co-Author Publications}
\cvitem{2026}{Dong, X., F. Massonnet, Y. Nie, B. Richaud, Y. Gao, L. Hoffman, Y. Wang, and Q. Yang. Incorporating subsurface oceanic variables improves seasonal Antarctic sea ice prediction with neural networks. \textit{Journal of Geophysical Research: Atmospheres}, 131, e2025JD045470. DOI: 10.1029/2025JD045470}
\cvitem{2025}{Giglio, D., J. Sala, J. Gilson, L. Hoffman, and B. Kawzenuk. Adaptive sampling of the upper ocean by autonomous floats during atmospheric river precipitation. \textit{Geophysical Research Letters}, 52, e2025GL117069. DOI: 10.1029/2025GL117069}
\cvitem{2020}{Ngo, T.T., L. Hoffman, G. Hoople, W. Trevena, U. Shakya, and G. Barr. Surface morphology and drug loading characterization of 3D-printed methacrylate-based polymer facilitated by supercritical carbon dioxide. \textit{The Journal of Supercritical Fluids}, 160, 104786. DOI: 10.1016/j.supflu.2020.104786}

\section{Patents}
\cvitem{2017}{McMahon, E., and L. Hoffman. High torque wind turbine blade, turbine, and associated systems and methods. U.S. Patent No. 9,797,370. U.S. Patent and Trademark Office. \url{https://patents.google.com/patent/US9797370B1/en}.}

\section{Invited Talks}
\cvitem{2025}{Earth Science Information Partners (ESIP) Machine Learning Cluster Meeting. Virtual. Oral presentation.}
\cvitem{2024}{Interagency Arctic Research Policy Committee (IARPC) Modelers' Community of Practice March 2024 Meeting -- Combining Modeling and Machine Learning Approaches to Understanding the Arctic Earth System. Virtual. Oral presentation.}
\cvitem{2023}{AI for the Study of Environmental Risk (AI4ER) Seminar Series, University of Cambridge. Cambridge, UK.}
\cvitem{2022}{American Meteorological Society Collective Madison Meeting, 17th Conference on Polar Meteorology and Oceanography. Madison, WI. Oral presentation.}
\cvitem{2022}{NOAA Science Seminar Series. Virtual. Oral presentation.}

\section{Conference Presentations, Workshops, and Seminars}
\cvitem{2025}{Kavli Institute for Theoretical Physics (KITP), The Physics of Changing Polar Climate. Santa Barbara, CA. Oral presentation.}
\cvitem{2025}{American Meteorological Society Denver Summit, 18th Conference on Polar Meteorology and Oceanography. Denver, CO. Oral presentation.}
\cvitem{2025}{Earth and Climate (ELIC) Seminar Series, Université catholique de Louvain. Louvain-la-Neuve, Belgium. Seminar.}
\cvitem{2024}{Physical Oceanography Dissertation Symposium (PODS). Lihue, HI. Oral presentation.}
\cvitem{2024}{Workshop on the Role of Sea Ice and its Variability in the Climate System, International Centre for Theoretical Physics (ICTP). Trieste, Italy. Poster.}
\cvitem{2024}{Earth and Climate (ELIC) Seminar Series, Université catholique de Louvain. Louvain-la-Neuve, Belgium. Seminar.}
\cvitem{2023}{Mentoring Physical Oceanography Women to Increase Retention (MPOWIR) Pattullo Conference. Warrenton, VA.}
\cvitem{2023}{54th International Liège Colloquium on Ocean Dynamics: Machine learning and data analysis in oceanography. Liège, Belgium. Poster.}
\cvitem{2023}{SeaSAR 2023. Svalbard, Norway.}
\cvitem{2023}{Atmospheric and Ocean Sciences Forum, University of Colorado Boulder. Boulder, CO. Seminar.}
\cvitem{2023}{Scientific Machine Learning Symposium. San Diego, CA. Poster.}
\cvitem{2023}{ECCO Annual Meeting. Pasadena, CA. Oral presentation.}
\cvitem{2022}{9th Annual FIRO Workshop. San Diego, CA. Poster.}
\cvitem{2022}{Observing, Modeling, and Understanding the Circulation of the Arctic Ocean and Sub-Arctic Seas Workshop. Seattle, WA. Poster.}
\cvitem{2022}{Center for Western Weather and Water Extremes (CW3E) Annual Meeting. San Diego, CA. Poster.}
\cvitem{2022}{Ocean Salinity Conference. New York, NY. Oral presentation.}
\cvitem{2022}{UCSD Jacobs School of Engineering Research Expo. San Diego, CA. Poster.}
\cvitem{2022}{Ocean Sciences Meeting. Virtual. Oral presentation.}
\cvitem{2021}{American Meteorological Society 101st Annual Meeting. Virtual. Oral presentation.}
\cvitem{2020}{International Atmospheric Rivers Conference Sponsored Symposium. Virtual.}
\cvitem{2020}{Ocean Sciences Meeting. San Diego, CA.}
\cvitem{2016}{Institute of Electrical and Electronics Engineers Global Humanitarian Technology Conference. Seattle, WA. Poster.}

\section{Teaching Experience}
\cventry{Summer 2020 \& 2021}{Instructor, Climate Change and the Ocean}{University of California San Diego, Academic Connections}{San Diego, CA}{}{%
\begin{itemize}
    \item Course: Climate Change and the Ocean.
    \item Developed and implemented a 4-week online pre-college course for high school students, including research and communication skill-building through accessing, plotting, analyzing, and presenting real-world oceanographic data.
\end{itemize}}

\cventry{Sep. 2018--Dec. 2020}{Teaching Assistant}{University of California San Diego}{San Diego, CA}{}{%
\begin{itemize}
    \item Courses: introductory fluid mechanics (3 offerings), advanced fluid mechanics (5 offerings), and material and energy balances (1 offering).
    \item Led five hours of weekly discussion sessions with supplementary lectures and problem-solving support.
\end{itemize}}


\section{Guest Lectures}
\cvitem{Winter 2025}{Introduction to Machine Learning for Earth and Climate Science. \textit{Forecast, Prediction, and Projection in Climate Science}. UCLouvain.}
\cvitem{Winter 2023}{Introduction to Machine Learning for Earth and Climate Science. \textit{Analysis of Physical Oceanographic Data}. Scripps Institution of Oceanography.}

\section{Leadership, Service, and Mentorship}
\cventry{2025--Present}{Co-Lead}{OASIS Working Group}{}{}{%
\begin{itemize}
    \item Coordinate an international Antarctica InSync working group focused on AI-enabled approaches relevant to Antarctic sea ice observing-system design and recommendations.
\end{itemize}}
\cvitem{ORCAS}{Working Group Member.}
\cvitem{Peer Review}{Journal of Climate; Geophysical Research Letters; Journal of Geophysical Research: Atmospheres; npj; Environmental Modelling \& Software; Open Science Europe.}
\cvitem{Mentorship}{Antonio Martinez Soares, Master's Student, Université catholique de Louvain, Fall 2024--Spring 2025.}
\cvitem{}{Kayli Matsuyoshi, Undergraduate Student, Scripps Institution of Oceanography, Summer 2021.}
\cvitem{}{Aniruddh Varadarajan, Undergraduate Student, University of California San Diego, Fall 2020--Spring 2021.}

\section{Honors and Awards}
\cvitem{2021}{Interdisciplinary Research Award, UC San Diego.}
\cvitem{Mar.--Jun. 2020}{Teaching Assistant Commendation, ``For Extraordinary Impact as a Teaching Assistant.''}
\cvitem{2015--2016}{Achievement Rewards for College Scientists (ARCS), \$5,000.}
\cvitem{2011--2016}{Alcala Award, Merit Based Scholarship, \$20,000 annually.}
\cvitem{Fall 2015--Present}{Tau Beta Pi, Engineering Honor Society.}
\cvitem{Spring 2016--Present}{Pi Tau Sigma, Engineering Honor Society.}

\section{Activities}
\cvitem{Nov. 2022-- Dec. 2025}{MPOWIR Mentor Group.}
\cvitem{Jan.--Dec. 2021}{URGE Pod Member, Unlearning Racism in the Geosciences.}
\cvitem{Feb. 2013--Feb. 2015}{Engineers Without Borders, Fundraising Director.}

\section{Personal Interests}
\cvitem{}{Bouldering, backpacking, guitar, violin, yoga, and cats.}

\end{document}
